% =====================================================================
% Mid-term Report - Data Analysis with Python
% Topic 3: Clustering - HR Analytics Case Study
% Compile: pdflatex main.tex (run twice)
% Figures are loaded from ./figures/
% =====================================================================

\documentclass[12pt,a4paper]{article}

% ---------- Packages ----------
\usepackage[utf8]{inputenc}
\usepackage[T1]{fontenc}
\usepackage[english]{babel}
\usepackage[margin=2.5cm]{geometry}

\usepackage{amsmath,amssymb}
\usepackage{graphicx}
\usepackage{booktabs}
\usepackage{tabularx}
\usepackage{array}
\usepackage{float}
\usepackage{caption}
\usepackage{subcaption}
\usepackage{xcolor}
\usepackage{enumitem}
\usepackage{listings}
\usepackage{fancyhdr}
\usepackage{setspace}
\usepackage[hidelinks]{hyperref}

% ---------- Document formatting ----------
\setstretch{1.15}
\setlength{\parskip}{4pt}
\setlength{\parindent}{0pt}

% ---------- Header / footer ----------
\pagestyle{fancy}
\fancyhf{}
\lhead{\small Data Analysis with Python -- Mid-term Report}
\rhead{\small Class AI67C}
\cfoot{\small\thepage}
\renewcommand{\headrulewidth}{0.4pt}

% ---------- Helper macros ----------

% Include a figure if it exists, otherwise show a placeholder
\newcommand{\safefig}[3][0.85\textwidth]{%
  \IfFileExists{figures/#2}{%
    \includegraphics[width=#1]{figures/#2}%
  }{%
    \fbox{\parbox[c][4cm][c]{0.7\textwidth}{\centering\color{red}%
    Missing figure: \texttt{figures/\detokenize{#2}}\\[2pt]#3}}%
  }%
}

% ---------- Code listing style ----------
\lstset{
  basicstyle=\ttfamily\footnotesize,
  keywordstyle=\color{blue!70!black},
  commentstyle=\color{green!40!black},
  stringstyle=\color{orange!70!black},
  breaklines=true,
  frame=single,
  language=Python,
  numbers=left,
  numberstyle=\tiny\color{gray},
  showstringspaces=false
}

% =====================================================================
\begin{document}

% =====================================================================
%                         COVER PAGE
% =====================================================================

\begin{titlepage}
    \centering

    % Compact spacing for the cover page
    \setlength{\parskip}{0pt}
    \setstretch{1.0}

    % University name
    {\Large\bfseries NATIONAL ECONOMICS UNIVERSITY\par}

    \vspace{0.15cm}

    {\large\bfseries COLLEGE OF TECHNOLOGY\par}

    \vspace{0.3cm}

    % Decorative line
    {\large ---------- $\ast$ ----------\par}

    % University logo in the middle space
    \vspace{0.45cm}

    \IfFileExists{figures/neu_logo.png}{
        \includegraphics[width=0.19\textwidth]{figures/neu_logo.png}
    }{
        \fbox{
            \parbox[c][2.5cm][c]{0.19\textwidth}{
                \centering
                \color{red}
                Missing logo\\
                neu\_logo.png
            }
        }
    }

    % Assignment title
    \vspace{0.65cm}

    {\LARGE\bfseries MID-TERM ASSIGNMENT\par}

    \vspace{0.25cm}

    {\large COURSE: Data Analysis with Python\par}

    % Topic
    \vspace{0.65cm}

    {\large\bfseries
    TOPIC: CLUSTERING EMPLOYEES BY WORK BEHAVIOUR\\[3pt]
    AND SATISFACTION -- HR ANALYTICS CASE STUDY
    \par}

    % Advisor and class
    \vspace{0.75cm}

    {\normalsize\bfseries Academic advisor: Mr. Tran Duc Minh\par}

    \vspace{0.15cm}

    {\normalsize\bfseries Class: AI67C\par}

    % Members table
    \vspace{0.45cm}

    \begin{tabular}{p{5.5cm} p{3.5cm}}
        \toprule
        \textbf{Members} & \textbf{Student ID}\\
        \midrule
        Dang Hoang Minh TU & 11258574\\
        Tran Nguyen Khoi & 11253986\\
        Nguyen Hong Nhat & 11256438\\
        \bottomrule
    \end{tabular}

    % Academic year at the bottom
    \vfill

    {\large\bfseries 2026--2027\par}

\end{titlepage}

% =====================================================================
%                         ABSTRACT
% =====================================================================

\pagenumbering{roman}

\begin{abstract}
This report applies unsupervised learning to the \emph{HR Analytics Case Study}
dataset, which describes about 4{,}410 employees of a company through
demographic, job, survey and attendance (swipe-in/swipe-out) records.

The goal is to discover naturally occurring employee segments without using
any pre-defined label. We integrate the relational files, clean and engineer
features (including average daily working hours derived from attendance logs),
and compare three clustering algorithms --- K-Means, Hierarchical (Ward)
clustering and DBSCAN --- using the Silhouette score, the Elbow method,
Davies--Bouldin and Calinski--Harabasz indices.

The resulting segments are profiled and cross-checked against the
\texttt{Attrition} column (used only for external validation) to derive
practical HR recommendations.
\end{abstract}

% =====================================================================
%                         TABLE OF CONTENTS
% =====================================================================

\tableofcontents

\clearpage

% =====================================================================
%                         MAIN REPORT
% =====================================================================

\pagenumbering{arabic}

% Add your report chapters below this line.

% Example:
% \section{Introduction}
% Your content here.

% =====================================================================

\section{Initial Data Inspection}

Before data preprocessing, an initial inspection was conducted
to understand the structure and quality of all five datasets.
The inspection focused on the consistency of the files, the
relationships between them, file encoding, data types,
missing values, and duplicate records.

\subsection{Dataset Overview}

The HR Analytics dataset consists of five relational files:
\texttt{general}, \texttt{employee}, \texttt{manager},
\texttt{in\_time}, and \texttt{out\_time}.
Each file contains 4,410 records.

Table~\ref{tab:dataset-overview} summarizes the basic
information of the five datasets.

\begin{table}[H]
\centering
\caption{Overview of the five datasets}
\label{tab:dataset-overview}
\small
\begin{tabular}{lrrrrr}
\toprule
\textbf{Criteria}
& \textbf{General}
& \textbf{Employee}
& \textbf{Manager}
& \textbf{In-time}
& \textbf{Out-time}\\
\midrule
Records & 4,410 & 4,410 & 4,410 & 4,410 & 4,410\\
Attributes & 24 & 4 & 3 & 262 & 262\\
Missing values & 28 & 83 & 0 & 109080 & 109080\\
Duplicate records & 0 & 0 & 0 & 0 & 0\\
\bottomrule
\end{tabular}
\end{table}

Across the five files, there are 555 attributes in total.
However, the 22,050 records represent the sum of records
across files, not the number of unique employees, since the
same employees are represented in multiple relational files.

\subsection{Column Consistency}

The five files describe different information, so they cannot be
concatenated; they must be merged later. Consistency was therefore
checked in two ways. First, the numbers of rows and columns were
compared (Table~\ref{tab:dataset-overview}): every file has 4,410 rows,
while the number of columns differs by file (24, 4, 3, 262 and 262).
Second, the column names of \texttt{in\_time} and \texttt{out\_time},
which have the same structure (one column per calendar day), were compared.
They contain the same 262 column names in the same order, so the two
files can be processed in parallel when computing working hours.
No column name other than \texttt{EmployeeID} is shared between the
other files, so no naming conflict (e.g.\ \texttt{\_x},
\texttt{\_y} suffixes) will appear after merging.

\subsection{Relationships Between Files}

The common column of the five files is \texttt{EmployeeID}. In
\texttt{in\_time} and \texttt{out\_time} this column is stored in an
unnamed first column (\texttt{Unnamed: 0}) and is renamed to
\texttt{EmployeeID}. Table~\ref{tab:key-check} shows the key checks.

\begin{table}[H]
\centering
\caption{Primary key / foreign key check on \texttt{EmployeeID}}
\label{tab:key-check}
\small
\begin{tabular}{lrrrr}
\toprule
\textbf{File} & \textbf{Null keys} & \textbf{Duplicate keys}
& \textbf{IDs not in General} & \textbf{General IDs missing}\\
\midrule
General  & 0 & 0 & 0 & 0\\
Employee & 0 & 0 & 0 & 0\\
Manager  & 0 & 0 & 0 & 0\\
In-time  & 0 & 0 & 0 & 0\\
Out-time & 0 & 0 & 0 & 0\\
\bottomrule
\end{tabular}
\end{table}

\texttt{EmployeeID} is the \textbf{primary key} of the \texttt{general}
dataset (unique and non-null) and the \textbf{foreign key} in the other
four files, so the relationship is one-to-one. A test merge of
\texttt{general}, \texttt{employee} and \texttt{manager} on
\texttt{EmployeeID} returned 4,410 rows, i.e.\ no record was lost
or duplicated.

\subsection{File Encoding}

The encoding of each CSV file was tested by decoding the raw bytes.
All five files were successfully decoded as UTF-8 (no BOM, no decoding
error), so no special encoding argument is needed when calling
\texttt{pd.read\_csv()} and no font or character errors are expected.

\subsection{Data Types}

The data types of each dataset were examined using the
\texttt{info()} method.

\begin{table}[H]
\centering
\caption{Distribution of data types across datasets}
\label{tab:data-types}
\small
\begin{tabular}{lrrrr}
\toprule
\textbf{Dataset}
& \textbf{int64}
& \textbf{float64}
& \textbf{str}
& \textbf{Total attributes}\\
\midrule
General & 14 & 2 & 8 & 24\\
Employee & 1 & 3 & 0 & 4\\
Manager & 3 & 0 & 0 & 3\\
In-time & 1 & 12 & 249 & 262\\
Out-time & 1 & 12 & 249 & 262\\
\midrule
\textbf{Total} & \textbf{20} & \textbf{29}
& \textbf{506} & \textbf{555}\\
\bottomrule
\end{tabular}
\end{table}

The \texttt{general} dataset contains demographic and job-related
information. The \texttt{employee} dataset provides employee
satisfaction and work-life balance information, while the
\texttt{manager} dataset contains job involvement and
performance ratings.

The \texttt{in\_time} and \texttt{out\_time} datasets each
contain 262 attributes, including 249 string attributes,
12 floating-point attributes, and 1 integer attribute.
These two files contain attendance-related records and will
be examined further during feature engineering.

\subsection{Missing Values}

Missing values were identified in the \texttt{general} and
\texttt{employee} datasets.

\subsubsection{General Dataset}

The \texttt{general} dataset contains 28 missing values
across two attributes.

\begin{table}[H]
\centering
\caption{Missing values in the General dataset}
\label{tab:general-missing}
\begin{tabular}{lrr}
\toprule
\textbf{Attribute} & \textbf{Missing values}
& \textbf{Percentage}\\
\midrule
NumCompaniesWorked & 19 & 0.43\%\\
TotalWorkingYears & 9 & 0.20\%\\
\midrule
\textbf{Total} & \textbf{28} & \textbf{0.63\%}\\
\bottomrule
\end{tabular}
\end{table}

\subsubsection{Employee Dataset}

The \texttt{employee} dataset contains 83 missing values
across three attributes.

\begin{table}[H]
\centering
\caption{Missing values in the Employee dataset}
\label{tab:employee-missing}
\begin{tabular}{lrr}
\toprule
\textbf{Attribute} & \textbf{Missing values}
& \textbf{Percentage}\\
\midrule
EnvironmentSatisfaction & 25 & 0.57\%\\
JobSatisfaction & 20 & 0.45\%\\
WorkLifeBalance & 38 & 0.86\%\\
\midrule
\textbf{Total} & \textbf{83} & \textbf{1.88\%}\\
\bottomrule
\end{tabular}
\end{table}

\subsubsection{Manager Dataset}

The \texttt{manager} dataset contains three attributes:
\texttt{EmployeeID}, \texttt{JobInvolvement}, and
\texttt{PerformanceRating}. All 4,410 records are non-null,
indicating that no missing values were detected.

\subsubsection{Attendance Datasets}

The \texttt{in\_time} and \texttt{out\_time} datasets
each contain 262 attributes and 4,410 records.
The available inspection output does not provide the
non-null counts for individual columns. Therefore,
their exact numbers of missing values need to be
calculated separately before preprocessing.

\subsection{Duplicate Records}

Duplicate records were checked in all five datasets.
The results showed that no duplicate records were detected
in any of the files.

\subsection{Initial Data Quality Assessment}

Based on the initial inspection, the following observations
were identified:

\begin{itemize}
    \item All five datasets contain 4,410 records each,
    with a total of 555 attributes.

    \item The datasets contain numerical and string data types,
    requiring appropriate preprocessing techniques.

    \item The \texttt{general} dataset contains 28 missing values,
    while the \texttt{employee} dataset contains 83 missing values.

    \item The \texttt{manager} dataset contains no missing values.

    \item The exact number of missing values in the attendance
    datasets still needs to be determined.

    \item No duplicate records were detected in the five datasets.

    \item \texttt{EmployeeID} links all files (primary key in
    \texttt{general}, foreign key elsewhere), and the files can be
    read with UTF-8 encoding.
\end{itemize}

These findings provide an overview of the data structure
and initial data quality. The next step is to examine the
attendance records, handle missing values, and prepare
suitable features for clustering.
% =====================================================================
\section{Data Cleaning}

The inspection in the previous section showed missing values in the
\texttt{general}, \texttt{employee}, \texttt{in\_time} and
\texttt{out\_time} files, no duplicate records, and attendance dates stored
as strings. This section describes how each problem was handled to obtain the
final table \texttt{df\_clean}. The pipeline is implemented in the
accompanying notebook.

% ---------------------------------------------------------------------
\subsection{Merging the Tables}

The \texttt{general}, \texttt{employee} and \texttt{manager} files were joined
on \texttt{EmployeeID} with \texttt{validate=}\allowbreak\texttt{"one\_to\_one"}, which raises an
error if any key is repeated. The merged table keeps all 4,410 employees and
contains 28 columns. The two attendance files are processed separately
(Section~\ref{sec:attendance}) and joined at the end.

% ---------------------------------------------------------------------
\subsection{Removing Duplicates}

Duplicates were checked again after merging, because joins can silently
multiply rows. Two criteria were used: fully identical rows and repeated
\texttt{EmployeeID}. Both checks returned 0, so no record was removed and the
table keeps 4,410 employees.

% ---------------------------------------------------------------------
\subsection{Handling Missing Values}

Missing values were not treated in the same way everywhere. The decision
depends on \emph{why} a value is missing.

\begin{table}[H]
\centering
\caption{Missing-value handling strategy}
\label{tab:missing-strategy}
\small
\begin{tabularx}{\textwidth}{l l r X}
\toprule
\textbf{Dataset} & \textbf{Attribute} & \textbf{Missing} & \textbf{Strategy}\\
\midrule
General  & NumCompaniesWorked      & 19 & Median imputation\\
General  & TotalWorkingYears       & 9  & Median imputation\\
Employee & EnvironmentSatisfaction & 25 & Mode imputation\\
Employee & JobSatisfaction         & 20 & Mode imputation\\
Employee & WorkLifeBalance         & 38 & Mode imputation\\
Manager  & ---                     & 0  & No action\\
In/Out-time & all date columns     & 109,080 each & Not imputed: interpreted as absence and
                                   converted to features\\
\bottomrule
\end{tabularx}
\end{table}

\subsubsection{General and Employee datasets}

The General and Employee datasets contain 28 (0.63\%) and 83 (1.88\%) missing
values, respectively. Because these counts are small, all rows were retained
instead of being deleted, so that the other attributes of the affected
employees are not lost. The missing values were imputed with a statistic
chosen according to the nature of each variable.

\paragraph{Count variables: median.}
\texttt{NumCompaniesWorked} and \texttt{TotalWorkingYears} are non-negative
counts. Figure~\ref{fig:figure1} shows that both are right-skewed (skewness
1.12 and 1.03): most employees have low values, while a small number have very
high ones (up to 40 years of experience and up to 9 previous companies) that
form a long upper tail. These extreme values pull the mean above the median
($11.28$ vs.\ $10$ years; $2.69$ vs.\ $2$ companies), so the mean no longer
describes a typical employee, and for an integer count it also produces
non-integer values. The median is not affected by the tail, so it was used.

\begin{figure}[H]
\centering
\safefig{figure1.png}{Distribution of TotalWorkingYears and NumCompaniesWorked}
\caption{Distribution of \texttt{TotalWorkingYears} and
\texttt{NumCompaniesWorked}, with mean (dashed) and median (solid).}
\label{fig:figure1}
\end{figure}

\paragraph{Ordinal survey scores: mode.}
\texttt{EnvironmentSatisfaction}, \texttt{JobSatisfaction} and
\texttt{WorkLifeBalance} are ordinal scores from 1 to 4. Figure~\ref{fig:figure2}
shows that their means (2.72, 2.73 and 2.76) fall between two levels and do
not correspond to any valid answer. The mode (3, 4 and 3) is the level chosen
most often, is always a valid score, and keeps the discrete nature of the
variable, so it was used instead of the mean.

\begin{figure}[H]
\centering
\safefig{figure2.png}{Distribution of the three survey scores with mean and mode}
\caption{Distribution of the three survey scores, with mean (dashed) and
mode (solid).}
\label{fig:figure2}
\end{figure}

% ---------------------------------------------------------------------
\subsection{Processing Attendance Data}
\label{sec:attendance}

Each of \texttt{in\_time} and \texttt{out\_time} has 261 date columns and
109,080 missing cells. These are \emph{not} random gaps. A cell is empty when
the employee did not swipe on that day, that is, on a public holiday, on leave
or when absent. Filling them with a mean or median would invent working hours
that never happened. The cells were therefore handled in five steps:

\begin{enumerate}[label=\textbf{Step \arabic*.}, leftmargin=*]
  \item \textbf{Convert to datetime.} All 261 date columns were parsed
  with \texttt{pd.to\_datetime}.

  \item \textbf{Remove non-working days.} Days on which no employee has a
  swipe-in and swipe-out are public holidays or closed days and carry no
  information. 12 such columns (the float columns in
  Table~\ref{tab:data-types}) were dropped. They account for
  $12 \times 4{,}410 = 52{,}920$ of the missing cells in each file, and the
  remaining $109{,}080 - 52{,}920 = 56{,}160$ cells are individual absences.
  The number of working days is therefore $N_{\text{work}} = 249$.

  \item \textbf{Validate records.} A day is valid only if both swipe-in
  and swipe-out exist and the duration lies in $(0, 24]$ hours. The check
  found 0 days with only one of the two swipes and 0 days with an invalid
  duration, so no record had to be discarded.

  \item \textbf{Aggregate per employee.} The remaining missing days are
  skipped (not filled) when computing the features below.

  \item \textbf{Join.} The three features are merged into the main table
  on \texttt{EmployeeID} (one-to-one).
\end{enumerate}

\begin{align}
  h_{i,d} &= t^{\text{out}}_{i,d} - t^{\text{in}}_{i,d}
  \quad\text{(hours, valid days only)},\\
  \text{AvgWorkHours}_i &= \frac{1}{|D_i|}\sum_{d\in D_i} h_{i,d},\\
  \text{AbsentDays}_i &= N_{\text{work}} - |D_i|,\qquad
  \text{AbsentRate}_i = \frac{\text{AbsentDays}_i}{N_{\text{work}}},
\end{align}

where $D_i$ is the set of valid days of employee $i$. Every employee has at
least one valid day, so no median fallback was needed for these features.

\begin{table}[H]
\centering
\caption{Summary of the engineered attendance features}
\label{tab:att-features}
\small
\begin{tabular}{lrrrrr}
\toprule
\textbf{Feature} & \textbf{Mean} & \textbf{Std} & \textbf{Min}
& \textbf{Median} & \textbf{Max}\\
\midrule
AvgWorkHours & 7.70  & 1.34 & 5.95  & 7.41  & 11.03\\
AbsentDays   & 12.74 & 5.50 & 1     & 13    & 24\\
AbsentRate   & 0.051 & 0.022 & 0.004 & 0.052 & 0.096\\
\bottomrule
\end{tabular}
\end{table}

After this step the dataset contains \textbf{no missing values}
(Table~\ref{tab:cleaning-summary}).

% ---------------------------------------------------------------------
\subsection{Standardizing Data Types}

\begin{table}[H]
\centering
\caption{Data type conversions}
\label{tab:dtype-conversion}
\small
\begin{tabularx}{\textwidth}{>{\raggedright\arraybackslash}X l l >{\raggedright\arraybackslash}X}
\toprule
\textbf{Attributes} & \textbf{Before} & \textbf{After} & \textbf{Reason}\\
\midrule
261 date columns of in/out-time & string & datetime & Needed to compute working hours\\
Survey scores, \texttt{NumCompaniesWorked}, \texttt{TotalWorkingYears}
  & float64 & int64 & Imputed values are whole numbers\\
\texttt{BusinessTravel}, \texttt{Department}, \texttt{EducationField},
  \texttt{Gender}, \texttt{JobRole}, \texttt{MaritalStatus}
  & string & category & Memory and correct encoding\\
\texttt{Attrition} & string (Yes/No) & int64 (1/0) & External validation only\\
\bottomrule
\end{tabularx}
\end{table}

Three constant attributes (\texttt{EmployeeCount}, \texttt{Over18},
\texttt{StandardHours}) have a single value for all employees. They cannot
separate clusters and were dropped.

A logical consistency check compared \texttt{YearsAtCompany} with
\texttt{TotalWorkingYears}. Two employees (0.05\%) have more years at the
company than total working years. This is negligible and these records were
kept unchanged.

% ---------------------------------------------------------------------
\subsection{Summary of Data Cleaning}

\begin{table}[H]
\centering
\caption{Dataset before and after cleaning}
\label{tab:cleaning-summary}
\begin{tabular}{lrr}
\toprule
\textbf{Criteria} & \textbf{Before} & \textbf{After}\\
\midrule
Employees (rows)                      & 4,410 & 4,410\\
Missing values (general + survey)     & 111   & 0\\
Missing values (attendance, per file) & 109,080 & 0 (converted to features)\\
Duplicate records                     & 0     & 0\\
Constant attributes                   & 3     & 0\\
Attributes (final table)              & 555 (5 files) & 29\\
\bottomrule
\end{tabular}
\end{table}

The 29 attributes of \texttt{df\_clean} are: 21 attributes of
\texttt{general} (24 minus 3 constant ones), 3 survey scores from
\texttt{employee}, 2 ratings (\texttt{JobInvolvement},
\texttt{PerformanceRating}) from \texttt{manager}, and 3 attendance
features. The cleaned table has one row per employee and is ready for
exploratory data analysis and feature selection.

\begin{lstlisting}[caption={Core of the attendance cleaning step}]
in_dt  = in_df.apply(pd.to_datetime, errors="coerce")
out_dt = out_df.apply(pd.to_datetime, errors="coerce")

# drop public holidays: days with no swipe for every employee
holiday = in_dt.columns[in_dt.isna().all() & out_dt.isna().all()]
in_dt, out_dt = in_dt.drop(columns=holiday), out_dt.drop(columns=holiday)

# keep only days with both in and out
half = in_dt.isna() ^ out_dt.isna()
in_dt, out_dt = in_dt.mask(half), out_dt.mask(half)

hours = (out_dt - in_dt).apply(lambda s: s.dt.total_seconds() / 3600)
hours = hours.mask((hours <= 0) | (hours > 24))

att = pd.DataFrame({
    "AvgWorkHours": hours.mean(axis=1),          # NaN days skipped
    "AbsentDays":   hours.isna().sum(axis=1),
})
att["AbsentRate"] = att["AbsentDays"] / hours.shape[1]
\end{lstlisting}

% =====================================================================
%                 LIST OF FIGURES AND LIST OF TABLES (END)
% =====================================================================

\clearpage
\listoffigures

\listoftables

\end{document}